\documentclass{../../../unipd-notes}

\unipdsetup{
  course = {Fondamenti di Ingegneria Informatica},
  short-course = {Fondamenti di Ingegneria},
  document-type = {Appunti delle lezioni}
}

\begin{document}
\chapter{Algoritmi di ricerca e accumulazione}

\begin{algorithm}[htbp]
  \caption{Ricerca lineare}
  \label{alg:linear-search}
  \KwInput{Una sequenza $a_1,\dots,a_n$ e un valore $x$}
  \KwOutput{La posizione di $x$, oppure $\bot$}
  \For{$i\gets 1$ \KwTo $n$}{
    \If{$a_i=x$}{
      \Return{$i$}\;
    }
  }
  \Return{$\bot$}\;
\end{algorithm}

Nel caso peggiore, l'\cref{alg:linear-search} confronta la chiave con tutti gli $n$ elementi e richiede tempo $\Theta(n)$.

\begin{algorithm}[htbp]
  \caption{Checksum additivo a otto bit}
  \label{alg:additive-checksum}
  \KwInput{Un frame di byte $b_1,\dots,b_n$}
  \KwOutput{Il byte di controllo $c$}
  $s\gets 0$\;
  \For{$i\gets 1$ \KwTo $n$}{
    $s\gets (s+b_i)\bmod 256$\;
  }
  \Return{$(-s)\bmod 256$}\;
\end{algorithm}
\end{document}
